\documentclass[11pt]{article}

\usepackage[margin=1in]{geometry}
\usepackage{hyperref}
\usepackage{xcolor}
\usepackage{listings}
\usepackage{enumitem}
\usepackage{parskip}

\definecolor{codebg}{RGB}{245,245,245}
\definecolor{codekeyword}{RGB}{0,0,150}
\definecolor{codecomment}{RGB}{110,110,110}
\definecolor{codestring}{RGB}{150,0,0}

\lstdefinestyle{pyblock}{
    backgroundcolor=\color{codebg},
    basicstyle=\ttfamily\small,
    keywordstyle=\color{codekeyword},
    commentstyle=\color{codecomment}\itshape,
    stringstyle=\color{codestring},
    showstringspaces=false,
    breaklines=true,
    frame=single,
    framerule=0.3pt,
    xleftmargin=1em,
    xrightmargin=1em,
    tabsize=4,
    captionpos=b,
    language=Python,
}
\lstdefinestyle{xmlblock}{
    backgroundcolor=\color{codebg},
    basicstyle=\ttfamily\small,
    showstringspaces=false,
    breaklines=true,
    frame=single,
    framerule=0.3pt,
    xleftmargin=1em,
    xrightmargin=1em,
    tabsize=4,
    captionpos=b,
    language=XML,
}

\hypersetup{
    colorlinks=true,
    linkcolor=blue,
    urlcolor=blue,
    citecolor=blue,
}

\newcommand{\code}[1]{\texttt{#1}}

\title{Python design description format}
\author{}
\date{}

\begin{document}

\maketitle

\tableofcontents

\section{Introduction}

This document specifies a new \code{XactFlow} \code{Importer}: a way to describe a
design's architecture (which IPs it instantiates, how they're wired together over buses,
and any ad hoc signal connections) as a plain Python script, instead of hand-written
IP-XACT XML. Running the script produces an \code{ipxact.Design} object, which the rest
of \code{XactFlow} already knows what to do with: \code{elaborate()} resolves and checks
it against a \code{Library}, and \code{XactFlow-design} can serialize it back out to a
real IP-XACT \code{design} XML document.

This builds on the design description concept from \emph{IP-XACT-Based SoC Generation:
Tools and SoCMake Integration} (M.~Travaillard, B.~Denkinger, August 2026), specifically
its Section 4.3 and the \code{design2ipxact} transformation of Section 5.2. That transformation
scoped a full SoCMake-integrated pipeline (\code{sv2ipxact}, \code{design2ipxact},
\code{ipxact2sv}, a shared bus library, radiation-hardening hooks, CMake integration);
this document only covers the piece of it that becomes a standalone \code{XactFlow}
\code{Importer}, with no SoCMake or CMake dependency, consistent with every other
\code{XactFlow} plugin.

\section{Scope}

What this importer does:

\begin{enumerate}
    \item Execute a Python script that declares IP instances, bus connections, and ad hoc
    connections using the API below.
    \item Resolve each IP reference to a real \code{ipxact.Component} via a
    \code{Library} the script itself provides.
    \item Build the corresponding \code{ipxact.Design} object: one
    \code{ComponentInstance} per IP instance, one \code{Interconnection} per bus
    connection, one \code{AdHocConnection} per ad hoc connection.
\end{enumerate}

What it deliberately does not do, because \code{XactFlow} already does it elsewhere:

\begin{itemize}
    \item \textbf{Bus protocol compatibility checking.} \code{connect\_bus()} only
    matches interfaces by bus definition VLNV and interface mode (see below); it does not
    check signal-level compatibility, that's \code{xactflow.elaborate()}'s job, via its
    interconnection-compatibility SCR checks, run after this importer hands off the
    \code{Design}.
    \item \textbf{Address conflict checking, memory map validation.} Out of scope for the
    reasons in Section~\ref{sec:address-mapping}.
    \item \textbf{Anything else IP-XACT's semantic consistency rules cover.} This
    importer's only job is producing a syntactically well-formed \code{Design}; whether
    that \code{Design} is semantically sound, correct cross-references, compatible
    interfaces, valid addressing, sound hierarchy, and so on, is \code{SCR}'s job,
    unchanged from how every other \code{Design} (hand-written or otherwise imported) is
    checked today.
\end{itemize}

\section{The \texttt{Design} object}

\begin{lstlisting}[style=pyblock, caption={The \texttt{Design} API surface.}]
class Design:
    def __init__(self, vlnv: str, library: xactflow.Library) -> None: ...
    def add_ip(self, instance_name: str, vlnv: str, parameters: dict[str, object] | None = None, transforms: list[str] | None = None) -> IPInstance: ...
    def connect_bus(self, a: IPInstance | BusInterfaceRef, b: IPInstance | BusInterfaceRef, *, base_addr: int | None = None, size: int | None = None) -> None: ...
    def connect(self, source: PortRef, target: PortRef) -> None: ...
    def build(self) -> ipxact.Design: ...
\end{lstlisting}

\subsection{\texttt{Design(vlnv, library)}}

\begin{itemize}
    \item \code{vlnv}: a colon-separated \code{"vendor:library:name:version"} string
    identifying the design itself. Parsed eagerly; a malformed string (not exactly 4
    fields) raises \code{ValueError} at construction.
    \item \code{library}: an already-scanned \code{xactflow.Library}. The script builds
    this itself (\code{Library.scan(...)}) and passes it in, \code{Design} never
    constructs or owns a \code{Library} implicitly. This keeps the script fully
    self-contained.
\end{itemize}

\subsection{\texttt{add\_ip(instance\_name, vlnv, parameters=None, transforms=None) -> IPInstance}}

Declares one component instance and resolves it immediately, not deferred to
\code{build()}:

\begin{itemize}
    \item \code{instance\_name} must be unique within the design. A duplicate raises
    \code{ValueError} immediately, naming both the new call and the earlier one that
    already used it.
    \item \code{vlnv} is resolved against \code{library} right away. An unresolvable VLNV
    raises \code{ValueError} at the \code{add\_ip()} call site, not at the end of the
    script, so a mistake is reported on the line that made it.
    \item \code{parameters} maps a parameter's \textbf{name} (as declared on the resolved
    component) to an override value. This is IP-XACT's own \code{configurableElementValues}
    feature. For each entry, \code{add\_ip} looks up the matching \code{Parameter} by name
    on the resolved component:
    \begin{itemize}
        \item Unknown name: \code{ValueError}, naming the instance, the bad key, and the
        parameter names that \emph{are} available on that component.
        \item Name matches a parameter with no \code{parameterId}: \code{ValueError}.
        Such a parameter cannot be referenced for override at all, per the standard; this
        is a defect in the referenced \emph{component} file, not something this importer
        can paper over (see Section~\ref{sec:relationship-sv} for why this matters in
        practice).
        \item Otherwise, the override is recorded against that parameter's
        \code{parameterId}.
    \end{itemize}
    \item \code{transforms} names zero or more transformations to apply to this
    instance, radiation hardening being one example. This importer does not apply them.
    It only records the names, in order, as a vendor extension on the resulting component
    instance.
    A separate tool, a transformer, does the work later. It takes the
    \code{ipxact.Design}, applies the transformations its instances ask for, and returns
    the modified \code{ipxact.Design}, ready for elaboration. A transformer can be kept
    private, and can be turned on or off without changing the design description.
\end{itemize}

Returns an \code{IPInstance} handle (see below) used to build bus and ad hoc connections
against this instance.

\subsubsection{\texttt{IPInstance}}

A handle returned by \code{add\_ip()}, carrying the instance name and the resolved
\code{ipxact.Component}.

\begin{lstlisting}[style=pyblock, caption={The \texttt{IPInstance} handle.}]
class IPInstance:
    name: str
    component: ipxact.Component

    def port(self, name: str) -> PortRef: ...
    def bus(self, name: str) -> BusInterfaceRef: ...
\end{lstlisting}

\begin{itemize}
    \item \code{.port(name)} looks up \code{name} in \code{component.model.ports};
    unknown name raises \code{ValueError} listing the ports that do exist. Used to build
    ad hoc connections.
    \item \code{.bus(name)} looks up \code{name} in \code{component.bus\_interfaces};
    unknown name raises \code{ValueError} the same way. Used only when
    \code{connect\_bus()} needs help picking the right bus (see below).
    \item Attribute access (\code{instance.some\_port}) is shorthand for
    \code{instance.port("some\_port")}, for the common case of wiring up ad hoc
    connections without the extra call. It raises the same \code{ValueError} as
    \code{.port()} for an unknown name. Note that \code{name}, \code{component},
    \code{port}, and \code{bus} are real attributes/methods of \code{IPInstance} itself,
    so a port that happens to share one of those names cannot be reached this way; the
    explicit \code{.port(...)} form always works instead.
\end{itemize}

\subsection{\texttt{connect\_bus(a, b, *, base\_addr=None, size=None)}}

Connects one bus interface on \code{a} to one bus interface on \code{b}. Each argument is
either a plain \code{IPInstance} (pick one of its bus interfaces automatically) or an
\code{IPInstance.bus(name)} (use that exact one). Here is how it picks:

\begin{enumerate}
    \item For each side, the candidates are: the exact interface given, if a
    \code{bus(name)} was passed; otherwise every bus interface on that instance's
    component that an earlier \code{connect\_bus()} call has not already used.
    \item A pair \code{(interface\_a, interface\_b)} matches if they share the same bus
    type, and their \code{mode}s are opposite. The same bus type means the same bus VLNV
    (\code{interface\_a.bus\_type.vlnv == interface\_b.bus\_type.vlnv}) and the same
    parameter values set on it (\code{bus\_type.config\_element\_values}), compared after
    evaluation. This first version only matches \code{INITIATOR} with
    \code{TARGET}; \code{SYSTEM}, \code{MIRRORED\_*}, and \code{MONITOR} modes are real
    IP-XACT concepts but are not matched automatically here, connecting through one needs
    an exact, named call instead. There is no need to say which side is the initiator and
    which is the target: each side's mode is already set on its own component, so this
    importer reads it instead of asking the script to repeat it.
    \item If exactly one matching pair exists, it is used.
    \item If more than one exists, \code{connect\_bus} raises \code{ValueError} listing
    every candidate interface name on both sides. The caller must then say exactly which
    one to use, by passing \code{a.bus("...")}/\code{b.bus("...")} for at least one side.
    A multi-port interconnect (several identical target interfaces, one per attached
    peripheral) will usually hit this case; it is expected, not an edge case.
    \item If none matches, \code{ValueError} explains why (no shared bus definition, no
    opposite mode available, or every candidate already used) and lists what bus
    interfaces are still free on each side.
\end{enumerate}

On success, one \code{ipxact.Interconnection} is added to the design, with
\code{active\_interface} set from \code{a}'s side and \code{other\_active\_interfaces=[<b's
side>]}. Its name always follows the same pattern, \code{f"\{a.name\}\_to\_\{b.name\}"} (a
number is added at the end if that name is already taken, e.g.\ a second link between the
same two instances). Both interfaces are now marked used, so a later
\code{connect\_bus()} call will not reuse them.

\subsubsection{Address mapping}
\label{sec:address-mapping}

\code{base\_addr}/\code{size} are a convenience for the common case: a target sitting at
some offset on an addressable bus. IP-XACT actually handles this through
\code{addressSpace}/\code{segment} on the initiator side, resolved through a
\code{designConfiguration}, but \code{XactFlow} doesn't elaborate that yet (see the main
\code{XactFlow} README's Known Limitations). Doing this the real way is future work, once
that support exists. Until then, when either is given,
they're recorded as a vendor extension on the \emph{target} side of the connection, the
same way IP-XACT represents anything the standard doesn't cover and isn't modeled here
yet. Concretely, something to the effect of:

\begin{lstlisting}[style=xmlblock, caption={Placeholder address-mapping vendor extension.}]
<xactflow:addressMapping baseAddr="0x40010000" size="0x1000"/>
\end{lstlisting}

is appended to that interface's \code{vendor\_extensions}. This is a placeholder, not a
standards-compliant addressing mechanism; anything reading it back needs to know to look
for it specifically. It should be revisited (and likely replaced) once \code{XactFlow}
elaborates \code{DesignConfiguration}/\code{addressSpace} for real.

\subsection{\texttt{connect(source, target)}}

An ad hoc, signal-level connection between exactly two ports, each a \code{PortRef}
(from \code{.port(name)} or the attribute shorthand). Unlike \code{connect\_bus}, this
validates immediately rather than deferring to \code{elaborate()}, since the check is
cheap and both ports are already in hand:

\begin{itemize}
    \item Both ports must exist (already guaranteed, \code{PortRef} can only be
    constructed from a real port lookup).
    \item Directions must be compatible: one \code{out} and one \code{in}, or either side
    \code{inout}. Two \code{out}s or two \code{in}s raise \code{ValueError} naming both
    ports and their directions.
\end{itemize}

On success, appends one \code{ipxact.AdHocConnection} with a single
\code{InternalPortReference} pair (\code{component\_instance\_ref}/\code{port\_ref} for
each side), named
\code{f"\{source.instance.name\}\_\{source.name\}\_to\_\{target.instance.name\}\_\{target.name\}"}.
Bit-slice or struct-sub-port ad hoc connections (\code{SubPortReference}/\code{PartSelect}
in \code{ipxact-compiler}'s model) are out of scope for this first version;
\code{connect()} only wires whole ports to whole ports.

\subsection{\texttt{build() -> ipxact.Design}}

Returns the finished \code{ipxact.Design}: every \code{add\_ip} as a
\code{ComponentInstance} (in call order), every \code{connect\_bus} as an
\code{Interconnection}, every \code{connect} as an \code{AdHocConnection}. It does not
check that every bus interface got connected: an unconnected optional interface (e.g.\ an
unused debug port) is a legitimate, common design, not an error. Completeness is
\code{SCR}'s concern, if and when a rule for it exists, not this importer's.

Scripts don't call \code{build()} themselves. The importer does, after running the
script, see below.

\section{How the importer runs a script}

\begin{lstlisting}[style=pyblock, caption={The \texttt{Importer} entry point.}]
class ArchiDescImporter(xactflow.Importer):
    def import_(self, source: Path, **options) -> ipxact.Design: ...
\end{lstlisting}

It executes \code{source} as a plain Python file (\code{runpy.run\_path}), then looks for
a module-level variable named \code{design}, of type \code{Design}. Missing or
wrong-typed \code{design} raises \code{ValueError} naming what was expected. It then
calls \code{.build()} on it and returns the result. Nothing about \code{Library} setup is
the importer's business, that already happened inside the script when it constructed its
own \code{Design}.

\section{Worked example}

\begin{lstlisting}[style=pyblock, caption={A complete design description, elaborated in the same script.}]
from xactflow import Library, elaborate
from xactflow_archidesc import Design

library = Library.scan("path/to/ip/library")
design = Design("cern.ch:soc:example:0.1", library)

cpu   = design.add_ip("cpu",   "openhwgroup:core:ibex:1.0", transforms=["tmrg"])
uart  = design.add_ip("uart",  "cern.ch:ip:uart:1.0", parameters={"BAUDRATE": 115200})
timer = design.add_ip("timer", "pulp:ip:timer:1.0")
debug = design.add_ip("debug", "cern.ch:ip:debug_module:1.0")
bus   = design.add_ip("apb_bus", "cern.socgen:apb_rt:0.1.0")

design.connect_bus(cpu, bus)
design.connect_bus(bus, uart,  base_addr=0x4001_0000, size=0x1000)
design.connect_bus(bus, timer, base_addr=0x4002_0000, size=0x1000)

design.connect(debug.debug_req_o, cpu.debug_req_i)
design.connect(cpu.debug_resp_o, debug.debug_resp_i)

elaborated = elaborate(design.build(), library)
\end{lstlisting}

\code{design.build()} produces the \code{ipxact.Design}: 5 \code{ComponentInstance}s, 3
\code{Interconnection}s (\code{cpu\_to\_apb\_bus}, \code{apb\_bus\_to\_uart},
\code{apb\_bus\_to\_timer}, one per \code{connect\_bus} call), and 2
\code{AdHocConnection}s. Passing it straight to \code{elaborate()} with the same
\code{library} resolves and checks it exactly as it would a hand-written or
hand-exported design, all in this one script, with no separate importer step needed.
Running the same script through \code{ArchiDescImporter} instead (letting the CLI or
another tool drive it) does that same \code{build()} call automatically; elaboration is
then up to whatever calls the importer, not shown here.

\section{Requirement on imported components}
\label{sec:relationship-sv}

\code{add\_ip}'s \code{parameters} argument only works for a component whose parameters
carry a \code{parameterId}; that's what \code{config\_element\_values} keys by. Any
component supplied to this format, whether produced by \code{XactFlow-sv}, another
importer, or hand-written IP-XACT, needs \code{parameterId} set on each parameter meant
to be overridable this way. One with no \code{parameterId} cannot be referenced for
override at all, per the standard. This is a general requirement on the components this
format consumes, not something specific to one importer.

\section{Open questions, deferred rather than silently resolved}

Things this first version leaves open, on purpose:

\begin{itemize}
    \item \textbf{Hierarchical design description.} A sub-design instantiated as a
    component inside a bigger one is not addressed here; this format currently describes
    exactly one flat design. \code{XactFlow}'s own elaborator doesn't recurse into nested
    designs yet either (see its README's ``No multi-level design hierarchy''
    limitation), so there is nothing downstream to hand a hierarchical result to yet
    regardless.
    \item \textbf{Address mapping}, as above: a real system, not a vendor-extension
    placeholder, once \code{XactFlow} elaborates
    \code{DesignConfiguration}/\code{addressSpace}.
    \item \textbf{Protocol bridges} (e.g.\ an APB-to-OBI adapter between two incompatible
    bus segments): not addressed. \code{connect\_bus} either finds a matching pair or
    fails; nothing here automatically adds a bridge component. A bridge, if one exists as
    a real component, is just another \code{add\_ip} the script author wires in
    explicitly on both sides, sourced and placed by the user, not made automatically by
    the tool.
    \item \textbf{System/mirrored/monitor interface modes}: real IP-XACT concepts,
    deliberately not auto-matched by \code{connect\_bus} in this first version (see
    above).
    \item \textbf{Bus type parameter values}: \code{connect\_bus} compares the parameter
    values set on each side's bus type (a data width, for example), not just the bus VLNV.
    These values are raw expressions (\code{"32"} and \code{"0x20"} are equal, a value can
    refer to an instance parameter), so they have to be evaluated before they can be
    compared. This first version only compares the bus VLNV, until \code{XactFlow}
    evaluates parameter values well enough to do it.
\end{itemize}

\end{document}
